% 第6章 有序与对称性破缺（Order and broken symmetry）
\chapter{有序与对称性破缺}

低温下自发有序的出现是凝聚态物理的基本现象。铁磁体、反铁磁体、液晶与超导体都是
有序相，固态本身也是如此。所有这些现象共享一些基本性质与特征。例如，它们都以
这样的温度依赖为标志：某个相关物理性质在临界温度 $T_{\mathrm{c}}$ 上下表现出显著
差异。对每种相都可以定义一个序参量（order parameter）：$T>T_{\mathrm{c}}$ 时为零，
$T<T_{\mathrm{c}}$ 时非零，因此它充当系统是否有序的指示器。对铁磁性，序参量就是
磁化强度。本章将说明每种有序相都伴随着一个对称性破缺（broken symmetry）——这一
概念将在下一节解释。

\section{对称性破缺}

铁磁体中选定了一个唯一的方向，所有原子磁矩都沿它排列：它们全都选择了“向上”而非
“向下”。这其实是个相当惊人的效应，因为底层的物理方程并不区分“上”与“下”——微观
物理具有实验观察到的基态所不具备的对称性。要深入这一谜题，必须更细致地考虑
对称性。

一个有用的例子是欧拉压杆（Euler strut，见图 6.1）：一根竖直杆，底端固定于地面，
顶端均匀加载。载荷加重时杆被压缩；超过临界重量它就会弯曲。向左弯还是向右弯取决
于载荷放置得如何对称等细节；但即使在“载荷绝对居中”的理想化情形下，杆仍然必须
选择弯曲的方向。它危如累卵、如立刀锋，一个随机的热涨落也许就足以使它向某一边
弯曲。轻载杆所固有的左右对称性在弯曲时被打破——这称为对称性破缺。

\begin{figure}
\centering
\includegraphics[width=0.6\textwidth]{fig6_01.pdf}
\caption{（a）欧拉压杆。（b）压力增大时压杆被压缩但保持左右对称。（c）力超过临界值
后压杆弯曲，打破左右对称。（d）压杆也可以向相反方向打破对称。}
\end{figure}

凝聚态体系中观察到许多类似特征。驱动对称性破缺转变的参数可以是力（如外加压强），
但更常是温度。图 6.2 画出液体与固体中的原子。液体冷却时系统只有轻微收缩，仍保持
极高的对称度；但低于临界温度（熔点）后液体变为固体，对称性被打破。乍看令人惊讶——
固体画出来比液体“显得”更对称：固体中原子整齐排列而液体里原子乱作一团。关键的观察
是：液体中任何一点与其他任何点完全相同——把系统对时间平均，每个位置被原子造访的
频率一样；不存在原子排齐的独特方向或轴。总之，液体具有完全的平移与旋转对称性。
而在固体中这种高度对称几乎全部丧失：图 6.2 所画的固体仍有残余对称——它不再在任意
旋转下不变，但在四重旋转（$\pi/2$、$\pi$、$3\pi/2$、$2\pi$）下不变；不再在任意
平移下不变，但在晶格基矢的整数组合平移下不变。所以并非所有对称都丧失了，而是
液态的高对称被打破了。

\mnote{关于对称性的术语（Terminology concerning symmetry）：若描述系统的哈密顿量
在对称群诸元素所对应的变换下不变，则系统具有该对称性。群是元素的集合加上组合它
们的运算，遵循一组特定规则：群是"封闭"的，即元素的组合仍是群的成员；群组合满足
结合律；逆元与单位元均有良好定义。离散对称性指元素可数的对称群，如立方体的旋转
对称群；连续对称性则有不可数的连续多个元素，如球的旋转对称群。若把群的对称操作
施加于整个系统而系统不变，则系统具有整体对称性（global symmetry）；若对各空间
点施加不同的对称操作后哈密顿量仍不变，则为局域对称性（local symmetry）。超导体
的规范对称性就是这样一种局域对称性。}

\begin{figure}
\centering
\includegraphics[width=0.4\textwidth]{fig6_02.pdf}
\caption{液--固相变。左：高温态（统计平均）具有完全的平移与旋转对称性。右：系统在
临界温度 $T_{\mathrm{c}}$ 以下变为固体，这些对称性被打破。}
\end{figure}

铁磁体情形类似（图 6.3）。居里温度 $T_{\mathrm{C}}$ 以上的铁磁体具有完全的旋转
对称性：所有方向等价，每个磁矩可以指向任意方向。$T_{\mathrm{C}}$ 以下系统为所有
自旋“选定”了唯一方向——高温态较高的旋转对称被打破；较低温态的旋转对称降低了
（只允许绕磁化轴——图 6.3 中的向上方向——的旋转）。

\begin{figure}
\centering
\includegraphics[width=0.4\textwidth]{fig6_03.pdf}
\caption{顺磁体--铁磁体相变。左：高温态（统计平均）具有旋转对称性。右：系统在居里
温度 $T_{\mathrm{C}}$ 以下变为铁磁体，该对称性被打破。}
\end{figure}

值得注意的重要一点：对称性不可能逐渐改变——某个特定的对称性要么存在要么不存在。
因此相变是尖锐的，有序态与无序态之间有清晰的分界。低温下有序的出现可以由相当
一般的热力学考虑理解：自由能 $F=E-TS$ 联系着能量与熵。为使自由能极小，低温下
系统选择能量最低的基态——通常是有序的，即最小化 E；温度升高时，最大化 S 的态
对最小化 F 更重要，于是更倾向于无序态。

并非所有相变都涉及对称性变化。图 6.4 是水的相图：液--气两区之间的边界线终止于
一个临界点。因此可以“作弊”绕过尖锐相变——沿一条避免不连续变化的路径穿过相图
（图 6.4 中的路径 C）。临界温度（647 K）以上，气体与液体仅以密度区分；气--液
转变不涉及对称性变化。相反，固--液转变涉及对称性变化，因此熔化曲线没有临界点。

\begin{figure}
\centering
\includegraphics[width=0.5\textwidth]{fig6_04.pdf}
\caption{水的相图。路径 A 与 B 越过相边界，C 则否。路径 A 涉及对称性的变化，
B 与 C 则不。}
\end{figure}

对称性破缺转变性质中令人困惑的一点是：破缺对称的基态并不具有哈密顿量的对称性。
把铁磁体冷却通过转变温度时，系统必须选定一个特定方向让所有自旋指向——尽管底层
物理并未挑出任何特定方向。这怎么会发生？

理解这一点，考虑一个非磁的例子很有用：氨分子（$\mathrm{NH_3}$，见图 6.5）。三个
氢原子构成等边三角形的三个顶点（图 6.5(a)、(b) 为侧视），氮原子可以在三角平面
略上方或略下方。这种金字塔形状源于氮上的孤对电子与三条氮--氢键争抢空间；它给
分子一个偶极矩，同时也在空间中定义了特定方向、打破了底层哈密顿量的对称性——
因此这个金字塔不可能是系统的稳定态。

事实上可以设想在图 6.5(a) 与 (b) 两个金字塔态之间变换：让氮原子穿过平面中央，
把四面体翻转——就像伞被大风吹翻那样。若以 x 表示氮原子到平面的距离，可以设想
一条势能曲线 $V(x)$，在两个稳定组态处各有一个极小（图 6.5(c)）。该势的两个最低
能量本征态示于图 6.5(d) 与 (e)：分别为第一激发态 $\phi_1(x)$ 与基态 $\phi_0(x)$。

因此图 6.5(a)、(b) 所代表的组态不是系统的稳定组态，而是能量本征态的叠加——它们
是亚稳态。图 6.5(a) 的态由 $(\phi_0(x)-\phi_1(x))/\sqrt{2}$ 表示，图 6.5(b) 的态
由 $(\phi_0(x)+\phi_1(x))/\sqrt{2}$ 表示。若把系统制备在图 6.5(a) 的亚稳态，会
发现它在它与图 6.5(b) 的亚稳态之间振荡，频率由 $\phi_0(x)$ 与 $\phi_1(x)$ 的能量
差给出：氨中为 24 GHz，氨微波激射器（maser）用的正是这一跃迁。若把氮换成更重的
磷制成 $\mathrm{PH_3}$，频率下降约十倍。系统越重，图 6.5(a) 亚稳态的时间常数越长。

铁磁体中我们有重得多的“态”：不只是更重的原子，而是更多的原子——也许 $10^{23}$
个。所有自旋指向单一方向的系统基态其实并不是系统的真正定态，但它是亚稳态，稳定
时间超过宇宙年龄。要到达另一个态（所有自旋反转），必须同时翻转系统中每一个自旋
——极不可能。

\begin{figure}
\centering
\includegraphics[width=0.6\textwidth]{fig6_05.pdf}
\caption{（a）氨 $\mathrm{NH_3}$ 是四面体分子，可存在于两种形式之一，另一种示于
（b）。氮原子有两个稳定位置，因此可以定义该系统的势能 $V(x)$。（c）势能曲线。
（d）第一激发态波函数 $\phi_1(x)$。（e）基态波函数 $\phi_0(x)$。}
\end{figure}

\section{模型}

为了讨论对称性破缺的一些后果，方便的做法是考虑一些简单的磁性模型并尝试求解。

\subsection{铁磁性的朗道理论}

俄国物理学家朗道（Lev Landau）给出了一个简便的、直接产生相变的模型，它源于
非常一般的考虑。把磁化强度为 M 的铁磁体的自由能写成 M 的幂级数。由于“上”与
“下”没有能量差别，幂级数不能包含 M 的任何奇次幂。于是自由能 $F(M)$ 可写为
\bio{Lev D. Landau\\（1908--1968）}
\begin{equation*}
F(M) = F_0 + a(T)M^{2} + bM^{4}\tag{6.1}
\end{equation*}
$F_0$ 与 b 是常数（设 $b>0$），$a(T)$ 依赖温度。若让 $a(T)$ 在转变温度 $T_{\mathrm{C}}$
处变号，可以证明该系统给出恰当的相变。在转变附近感兴趣的区域写
$a(T)=a_0(T-T_{\mathrm{C}})$（$a_0$ 为正常数）。求系统基态需使自由能极小，即解
$\partial F/\partial M=0$：
\begin{equation*}
2M\left[a_0(T-T_{\mathrm{C}}) + 2bM^{2}\right] = 0.\tag{6.2}
\end{equation*}
左边是两项之积，任一项可为零：
\begin{equation*}
M = 0 \quad\text{或}\quad M = \pm\left[\frac{a_0(T_{\mathrm{C}}-T)}{2b}\right]^{1/2}.\tag{6.3}
\end{equation*}
第二个条件只在 $T<T_{\mathrm{C}}$ 时有效，否则等于要给负数开平方。第一个条件在
$T_{\mathrm{C}}$ 上下都适用，但在 $T_{\mathrm{C}}$ 以下它只给出一个不稳定平衡点
（可由 $\partial^{2}F/\partial M^{2}$ 判断）。于是磁化强度遵循图 6.6(b) 的曲线：
$T\geq T_{\mathrm{C}}$ 时为零，$T<T_{\mathrm{C}}$ 时非零且正比于 $(T_{\mathrm{C}}-T)^{1/2}$。

\begin{figure}
\centering
\includegraphics[width=0.4\textwidth]{fig6_06.pdf}
\caption{（a）铁磁体的自由能 $F(M)$。（b）相应的磁化强度随温度的函数。}
\end{figure}

朗道研究相变的方法称为平均场理论（mean-field theory）：它假定所有自旋“感受”到
邻居产生的同一个平均交换场，该场正比于磁化强度。这与 5.1.1 节的外斯模型完全
相同。平均场理论是能构造出来描述多种相变的最简单理论，在各种情形给出相似结果；
在不同情形有不同名字，如合金有序--无序转变的 Bragg--Williams 理论。平均场理论
不能准确解释临界区，因为“样品各处相同”的假设在临界区恰恰最不成立。

平均场理论忽略关联与涨落，而它们在 $T_{\mathrm{C}}$ 附近极为重要。非常接近临界
温度时，序参量出现大涨落；临界区的特征是所有尺度上都有涨落。烧水时，水安静而不
起眼地变热，直到接近沸点才喧哗起来、冒泡翻腾。\fn{1}{这一现象的直观演示是临界乳光（critical opalescence）：透过临界点附近的气体看物体，影像模糊混浊。原因是临界点附近密度涨落强烈，引起折射率的大幅变化。}铁磁体在 $T_{\mathrm{C}}$ 附近发生同类效应。
刻画涨落的主导长度尺度是关联长度（correlation length）$\xi$：$T\to T_{\mathrm{C}}$
时 $\xi\to\infty$，即临界点处关联长度变为无穷。涨落的重要性也可从 $F(M)$ 的形式
看出：M=0 处的曲率 $(\partial^{2}F/\partial M^{2})_{M=0}$ 在 $T=T_{\mathrm{C}}$
趋于零，M=0 的极小变得宽广平坦（见图 6.6）。因此，尽管平均场理论易于求解，其
关于临界区的预言必须谨慎对待。
\bio{Ernst Ising\\（1900--1998）}
\mnote{第一章中自旋 $z$ 分量的算符写作 $\hat{S}_z$。现在还需要标记自旋所在的格点，
故把 z 写成上标、格点标号写成下标：$\hat{S}_i^z$ 即格点 $i$ 上自旋 $z$ 分量的
算符。}

\subsection{海森伯模型与伊辛模型}

理解固体磁行为的另一条路是考虑磁相互作用的特定微观模型。一个常被研究的模型是
最近邻海森伯模型：
\begin{equation*}
\hat{\mathcal{H}} = -\sum_{\langle ij\rangle}\mathsf{J}\,\mathbf{S}_i\cdot\mathbf{S}_j,\tag{6.4}
\end{equation*}
常数 $\mathsf{J}$ 是交换积分，求和号下的 $\langle ij\rangle$ 表示对最近邻求和。
自旋 $\mathbf{S}_i$ 被当作三维矢量——允许指向三维空间中任意方向。但求和可以在 1、
2 或 3 维晶格上进行。这里要区分自旋所在晶格的维数 d 与自旋自身的维数 D（一般称
D 为序参量的维数）。海森伯模型 D=3（自旋是三维矢量）；但这些自旋可以排在 1、2、
3 维（想排 4 维也行！）的晶格上，故 d 可为 1、2、3……

相关模型是伊辛模型（Ising model）：自旋只允许朝上或朝下，即只考虑自旋的 z 分量。
其哈密顿量为
\begin{equation*}
\hat{\mathcal{H}} = -\sum_{\langle ij\rangle}\mathsf{J}\,S_i^{z}S_j^{z}.\tag{6.5}
\end{equation*}
此处序参量维数 D=1（自旋只允许沿 $\pm z$）。但同样可以把这些一维自旋排在
$d=1,2,\ldots$ 的晶格上。

\subsection{一维伊辛模型（$D=1, d=1$）}

把伊辛自旋放在一维晶格上，我们将证明：没有相变。先考虑含 $N+1$ 个自旋的链
（须考虑 N 个相邻“键”）。哈密顿量为
\begin{equation*}
\hat{\mathcal{H}} = -2\mathsf{J}\sum_{i=1}^{N}S_i^{z}S_{i+1}^{z}\tag{6.6}
\end{equation*}
设 $\mathsf{J}>0$，基态为所有相邻自旋铁磁排列（图 6.7(a)）。
基态能量为 $-N\mathsf{J}/2$，因为所有 $S_i^{z}=+\frac{1}{2}$。现在加入一个“错误”
——单个缺陷（图 6.7(b)）。
这多付出能量 $E=\mathsf{J}$：把一个有利的相互作用（省 $\mathsf{J}/2$）变成不利的
（费 $\mathsf{J}/2$），能量变化为 $\mathsf{J}$。但熵增 $S=k_{\mathrm{B}}\ln N$——
缺陷可以放在 N 个位置中的任何一个。让链变得很长（$N\to\infty$）时，缺陷的能量
代价不变（$\mathsf{J}$），熵增益却变为无穷。性质由自由能 $F=E-TS$ 决定：只要
温度非零，熵的考虑就意味着缺陷的存在使自由能骤降至 $-\infty$——缺陷会自发形成，
实际上 $T>0$ 时不存在长程序。换句话说，临界温度为零。这一论证对一维晶格上的
所有模型都成立（一维中熵总是获胜），结论：一维中不可能有长程序。

\begin{figure}
\centering
\includegraphics[width=0.6\textwidth]{fig6_07.pdf}
\caption{一维伊辛模型。（a）基态为 $N+1$ 个自旋全部铁磁排列。（b）加入单个缺陷。}
\end{figure}

\subsection{二维伊辛模型（$D=1, d=2$）}

把伊辛自旋放在二维晶格上，则会出现向磁有序态的相变，其临界温度非零。因为制造
缺陷的能量代价与熵增益都按缺陷的周长标度（见图 6.8），能量与熵可以势均力敌，
谁也没有压倒性优势。该模型的严格求解是二十世纪统计物理的杰出成就之一
（Lars Onsager 于 1944 年解决），其解超出本书范围。这说明：即便表述简单的问题
也绝非容易求解。
\bio{Lars Onsager\\（1903--1976）}

\begin{figure}
\centering
\includegraphics[width=0.5\textwidth]{fig6_08.pdf}
\caption{二维伊辛模型。自旋向上标记为 +、向下标记为 $-$。向下自旋缺陷的能量与熵
都按缺陷边界的周长标度。}
\end{figure}

\section{对称性破缺的后果}

打破对称性有种种后果：
\begin{itemize}[itemsep=0.2em]
\item \textbf{相变}：系统在温度 $T_{\mathrm{c}}$ 处发生急剧的行为变化，称为改变了
相（如液体$\to$固体、顺磁体$\to$铁磁体等）。相变附近的区域称为临界区（见 6.4 节）。
\item \textbf{刚性}：打破对称后，系统对维持该破缺对称态有强烈的能量偏好；试图
改变系统打破对称的方式会遭到抵抗。因此晶体不易弯曲、铁磁体表现永磁体行为
（见 6.5 节）。
\item \textbf{激发}：$T=0$ 时系统完全有序；有限温度下，序参量中的激发削弱有序。
晶体中这些激发称为格波，量子化为声子；铁磁体中对应的模式称为自旋波，量子化为
磁振子（见 6.6 节）。
\item \textbf{缺陷}：若宏观样品相邻两部分以不同方式打破对称，其边界将含有缺陷：
如晶体中的位错或铁磁体中的畴壁（见 6.7 节）。
\end{itemize}

表 6.1 总结了几种破缺对称相的性质。每种相都有高温无序态与低温有序态、一个序参量、
一种刚性现象、一组代表相对有序态的波动型偏离的激发，以及与“不同空间区域以不同
方式打破对称”相联系的缺陷。

\begin{table}
\centering
\caption{破缺对称相的性质。$\rho_G$ 是与倒格矢 G 对应的空间频率的电荷密度傅里叶
分量；超导体中的复波函数为 $\psi=|\psi|e^{i\phi}$；电极化强度 P 是单位体积的
电偶极矩。}
\begin{small}
\begin{tabular}{lllllll}
\toprule
现象 & 高温相 & 低温相 & 序参量 & 激发 & 刚性现象 & 缺陷\\
\midrule
晶体 & 液体 & 固体 & $\rho_G$ & 声子 & 刚性 & 位错、晶界\\
铁磁体 & 顺磁体 & 铁磁体 & M & 磁振子 & 永磁性 & 畴壁\\
反铁磁体 & 顺磁体 & 反铁磁体 & M（子晶格上） & 磁振子 & （颇为微妙） & 畴壁\\
向列相液晶 & 液体 & 取向液体 & $S=\langle\frac{1}{2}(3\cos^{2}\theta-1)\rangle$ & 指向涨落 & 各种 & 向错、点缺陷\\
铁电体 & 无极性晶体 & 极性晶体 & P & 软模 & 铁电滞回 & 畴壁\\
超导体 & 正常金属 & 超导体 & $|\psi|e^{i\phi}$ & -- & 超导电性 & 磁通线\\
\bottomrule
\end{tabular}
\end{small}
\end{table}

晶体在衍射实验中给出布拉格峰。因此可用 $\rho_G$——与倒格矢 G 对应的空间频率的
电荷密度傅里叶分量——作序参量。反铁磁体与铁磁体类似，只是序参量是反铁磁体一套
子晶格上的磁化强度。超导体有复序参量 $\psi=|\psi|e^{i\phi}$：$|\psi|^{2}$ 代表
超导电子密度（对超流或玻色凝聚体，代表凝聚分数）。波函数的相位 $\phi$ 是被打破
的对称性：正常态中它在任何点可取任意值，超导态中它在整个样品上变得均匀。对向列相
液晶（见图 6.9），序参量为 $S=\langle\frac{1}{2}(3\cos^{2}\theta-1)\rangle$——
函数 $\frac{1}{2}(3\cos^{2}\theta-1)$ 的平均值，$\theta$ 是分子长轴与指向矢
（director）$\hat{\mathbf{n}}$（沿液晶分子平均取向的单位矢量）的夹角。各向同性
（液体）态对应 S=0；完全排齐对应 S=1。

\begin{figure}
\centering
\includegraphics[width=0.5\textwidth]{fig6_09.pdf}
\caption{向列相液晶含棒状分子。（a）低温下它们如固体般有序。（c）高温下处于液态。
（b）中间温度区为向列态：分子长轴排齐，尽管分子质心的位置并无序。}
\end{figure}

\begin{table}
\centering
\caption{各种模型的临界指数。}
\begin{small}
\begin{tabular}{lcccc}
\toprule
模型 & 平均场 & Ising & Ising & Heisenberg\\
\midrule
$D$ & 任意 & 1 & 1 & 3\\
$d$ & 任意 & 2 & 3 & 3\\
$\beta$ & $\frac{1}{2}$ & $\frac{1}{8}$ & 0.326 & 0.367\\
$\gamma$ & 1 & $\frac{7}{4}$ & 1.2378(6) & 1.388(3)\\
$\delta$ & 3 & 15 & 4.78 & 4.78\\
\bottomrule
\end{tabular}
\end{small}
\end{table}

\section{相变}

6.2 节给出了铁磁性的朗道理论——一种平均场理论：所有自旋感受同样的交换场，导出
转变以下磁化强度按 $(T_{\mathrm{C}}-T)^{1/2}$ 变化。真实体系中，磁化强度在转变
附近确实按 $(T_{\mathrm{C}}-T)^{\beta}$ 变化，但指数 $\beta$ 未必等于 $\frac{1}{2}$。
因此指数给出关于相变性质的重要信息。还可以定义若干类似指数，即临界指数（critical
exponents）。实验发现相变温度 $T_{\mathrm{C}}$ 附近
\begin{equation*}
\begin{array}{cccl}
\chi & \propto & (T-T_{\mathrm{C}})^{-\gamma} & \qquad T>T_{\mathrm{C}}\\
M & \propto & (T_{\mathrm{C}}-T)^{\beta} & \qquad T<T_{\mathrm{C}}\\
M & \propto & H^{1/\delta} & \qquad T=T_{\mathrm{C}}
\end{array}\tag{6.7}
\end{equation*}
$\beta$、$\gamma$、$\delta$ 是临界指数；各模型取值见表 6.2。

如前所述，平均场理论忽略关联与涨落，这意味着它无望给出 $T_{\mathrm{C}}$ 附近——
恰恰是想计算临界指数之处——的正确描述。然而事实证明：平均场理论是四维及以上
（$d\geq4$）体系的正确描述，并给出正确的临界指数。

尽管平均场理论无法完满解释四维以下体系的临界行为，但只要接受普适性（universality）
假说，只需考虑一小组代表性的理想统计模型即可计算任何物理体系的临界指数。该假说
基于这样的观察：临界指数看起来惊人地不依赖相变的类型——无论液--气、铁磁--顺磁、
超导--正常或任何其他。按该假说，对连续相变，临界指数只取决于：
（1）系统维数 d；（2）序参量维数 D（准确地说是序参量的对称性）；（3）力是短程
还是长程。
因此只需考察特定的普适类（universality classes）——即短程力与长程力下特定的
D、d 组合——在每类中挑最简单的模型计算指数即可。（上述假设适用于静态临界指数，
而非表征时间依赖性质的动态临界指数。）

现在列出有解已知的模型。精确可解的模型包括：
（1）$d=1$ 的所有情形——如前所述，不表现连续相变；
（2）$d\geq4$ 的所有情形——给出平均场解；
（3）长程相互作用的全部情形——给出平均场解；
（4）$d=2$、$D=1$——即二维伊辛模型（见表 6.2）；
（5）任意 $d$ 的 $D=\infty$ 情形——称为球模型（spherical model）。
遗憾的是，多数真实情形对应 $d=3$ 与短程相互作用——尚未精确求解。

1970 年代初发现了一种计算临界行为的方法，即便对不能精确求解的模型也适用：考虑
体积为 $L^{d}$ 的自旋块（$L$ 远小于关联长度 $\xi$）。若把该块的线性尺寸放大
（标度）$n$ 倍，哈密顿量将相应地被标度，可以考察序参量在这一变换下如何改变
（重正化）。记该变换为 $\tau$，则哈密顿量变换为 $\mathcal{H}'=\tau(\mathcal{H})$。
临界点处 $\xi\to\infty$，故必存在极限函数 $\mathcal{H}^{*}$ 使 $\mathcal{H}^{*}$
是 $\tau$ 的不动点，即 $\tau(\mathcal{H}^{*})=\mathcal{H}^{*}$。这一寻找标度变换
不动点的方法称为重正化群（renormalization group）方法——因为标度变换的集合构成
一个群。它在凝聚态物理中极为成功，并与量子场论中的重正化有深刻的联系。不过
重正化群的更详细处理超出本书范围。

团簇模型可用于描述铁磁体的顺磁区（$T>T_{\mathrm{C}}$）。它保留平均场方法的某些
简单性（从而保留预言能力），又尝试对涨落建模：精确计算单个有限尺寸自旋团簇的
组态，再把它与样品其余部分的平均场耦合。当然关联长度在 $T_{\mathrm{C}}$ 处发散，
故越接近转变团簇尺寸需要越大。若选择团簇尺寸不受限制、可以变化的团簇模型，可以
得到与数据极好的符合。\fn{2}{见 R.~V.~Chamberlin, Nature, \textbf{408}, 337 (2000)。}

\section{刚性}

打破对称涉及基态的选择：自旋该都朝上、朝下、朝左还是朝右？宏观样品的能量在其
整个体积内以同一种方式打破对称时最低。若试图让宏观样品的不同部分以不同方式打破
对称，将出现反映额外能量代价的力——由此产生广义刚性或刚度（stiffness）。这种
刚性与有序密切相关，以下对晶体的论证可以说明。

考虑晶格常数为 $a$ 的立方晶体，设想把某一平面的键拉伸，使 $a$ 变为 $a+u$，产生
应变 $\epsilon=u/a$（见图 6.10）。这产生等于 $G\epsilon$ 的应力（G 为弹性模量），
单个键上的力为 $G\epsilon a^{2}=Gau$，键中储存的能量为 $\frac{1}{2}Gau^{2}$。高温
（远高于德拜温度\fn{3}{德拜温度 $\Theta_{\mathrm{D}}$ 是最高被占据声子模的能量除以 $k_{\mathrm{B}}$——假定声子无色散。当然总有色散，但 $\Theta_{\mathrm{D}}$ 仍刻画了系统的典型声子能量。多数德拜温度在 100--1000 K 范围。$T\gg\Theta_{\mathrm{D}}$ 时声子模可作经典处理。}）下每个键的平均储存能量为 $\frac{1}{2}k_{\mathrm{B}}T$
（能量均分定理，见附录 E），于是
\begin{equation*}
\langle u^{2}\rangle = \frac{k_{\mathrm{B}}T}{Ga}\tag{6.8}
\end{equation*}
故 $\langle u^{2}\rangle\propto1/G$：G 为零则涨落发散，反之亦然。非零弹性模量
（即刚性）的存在与涨落的有限性（即晶体的稳定）相联系。

这就是说：既然有有序的晶体，就必有刚性的晶体。晶体不易形变（与液体相反），因为
存在正比于弹性模量与 $(\nabla u)^{2}$ 的弹性能（u 是晶格位移），且晶格把力传递到
宏观样品的另一端。固体因此是刚性的。

这些想法同样适用于铁磁体。铁磁体中自旋排齐，把它们相对彼此转动要付出能量——
这就是永磁现象。磁化强度非均匀时有正比于 $(\nabla\mathbf{M})^{2}$ 的交换代价
（见 4.2.7 节，尤其式 4.18）。

最后一个例子：超导体中波函数的相位 $\phi$ 是序参量。超导体中 $\phi$ 在样品上
均匀，存在正比于 $(\nabla\phi)^{2}$ 的“相位刚度”能量。（事实上，把相位沿样品扭转
会产生正比于 $\nabla\phi$ 的超流。）

\begin{figure}
\centering
\includegraphics[width=0.5\textwidth]{fig6_10.pdf}
\caption{晶格参数为 $a$ 的立方晶体被拉伸，产生应变 $\epsilon=u/a$。}
\end{figure}

\section{激发}

固体在 T=0 有序，尽管零点涨落意味着即便此时原子也并非纯静态。非零温度下，有序
被热激发的晶格振动破坏——它们量子化为声子。声子的行为由色散关系刻画，即角频率
$\omega$ 与波矢 q 之间的关系（等价地，能量 $\hbar\omega$ 与动量 $\hbar q$ 之间的
关系）。一维单原子链的例子示于图 6.11。关键特征：声学声子在 q=0 处 $\omega=0$，
故只要波长 $\lambda=2\pi/q$ 足够长，产生一个波矢 q 的声子所需能量小到可以忽略。
只要温度非零声学声子就能被热激发，是因为从系统基态（q=0、$\omega=0$）到最低
声学声子能级之间没有需要跨越的能隙——声子色散关系在 q=0 处没有能隙（与光学声子
相反）。三维晶体中声学声子给出低温 $T^{3}$ 热容。

\begin{figure}
\centering
\includegraphics[width=0.5\textwidth]{fig6_11.pdf}
\caption{质量为 m 的原子由弹簧常数 K 的键连接的一维单原子链的声子色散。这些声子
是声学的：没有能隙（q=0 处 $\omega=0$），小 q 时 $\omega\approx v_{\mathrm{s}}q$
（$v_{\mathrm{s}}$ 为声速）。色散是无质量的，因为 q=0 处 $\omega=0$（即能量
$E=\hbar\omega=0$ 于动量 $p=\hbar q=0$）。相对论性粒子的能量满足
$E^{2}=p^{2}c^{2}+m^{2}c^{4}$，零动量能量为 $mc^{2}$——正比于质量。}
\end{figure}

每当打破一个连续的整体对称（如从液体造出固体、从顺磁体造出铁磁体），就可以以
可忽略的小能量代价在序参量中产生长波长激发。这类激发称为戈德斯通模（Goldstone
modes，有时称戈德斯通玻色子）。因为不费能量，它们是“无质量的”。粒子物理中的
例子是光子——一个戈德斯通玻色子。超导体的情形相当不同：打破的是连续的局域对称，
不产生戈德斯通模——原因相当微妙，与希格斯机制有关。

铁磁体在 T=0 完全有序；非零温度下有序被自旋波破坏，自旋波量子化为磁振子
（magnons）。关键特征：只要波长足够长，产生一个磁振子所需能量小到可以忽略。
磁振子在铁磁体中扮演声子在固体中的角色，是系统的戈德斯通模。下面将证明：对各向
同性铁磁体，磁振子色散关系在 q=0 处没有能隙。

\subsection{磁振子}

本节推导各向同性铁磁体的磁振子色散关系。这是重要问题，故给出两种推导：第一种
半经典，第二种量子力学。

（1）从自旋波色散的半经典推导开始。回忆海森伯模型的哈密顿量
\begin{equation*}
\hat{\mathcal{H}} = -\sum_{\langle ij\rangle}\mathsf{J}\,\hat{\mathbf{S}}_i\cdot\hat{\mathbf{S}}_j\tag{6.9}
\end{equation*}
（即式 6.4）。一维链中每个自旋有两个邻居，哈密顿量约化为
\begin{equation*}
\hat{\mathcal{H}} = -2\mathsf{J}\sum_i\hat{\mathbf{S}}_i\cdot\hat{\mathbf{S}}_{i+1}.\tag{6.10}
\end{equation*}
用式 C.7 计算 $\langle\hat{\mathbf{S}}_j\rangle$ 的时间依赖：
\begin{align*}
\frac{\mathrm{d}\langle\hat{\mathbf{S}}_j\rangle}{\mathrm{d}t} &= \frac{1}{\mathrm{i}\hbar}\langle[\hat{\mathbf{S}}_j, \hat{\mathcal{H}}]\rangle\tag{6.11}\\
&= -\frac{2\mathsf{J}}{\mathrm{i}\hbar}\langle[\hat{\mathbf{S}}_j, \cdots+\hat{\mathbf{S}}_{j-1}\cdot\hat{\mathbf{S}}_j+\hat{\mathbf{S}}_j\cdot\hat{\mathbf{S}}_{j+1}+\cdots]\rangle\tag{6.12}\\
&= -\frac{2\mathsf{J}}{\mathrm{i}\hbar}\langle[\hat{\mathbf{S}}_j, \hat{\mathbf{S}}_{j-1}\cdot\hat{\mathbf{S}}_j]+[\hat{\mathbf{S}}_j, \hat{\mathbf{S}}_j\cdot\hat{\mathbf{S}}_{j+1}]\rangle\tag{6.13}\\
&= \frac{2\mathsf{J}}{\hbar}\langle\hat{\mathbf{S}}_j\times(\hat{\mathbf{S}}_{j-1}+\hat{\mathbf{S}}_{j+1})\rangle.\tag{6.14}
\end{align*}
现在把每个格点上的自旋当经典矢量处理。系统基态是所有自旋排齐（设沿 $z$ 轴）：
$S_j^{z}=S$、$S_j^{x}=S_j^{y}=0$。考虑对该态的小偏离：$S_j^{z}\approx S$、
$S_j^{x},S_j^{y}\ll S$，于是
\begin{align*}
\frac{\mathrm{d}S_j^{x}}{\mathrm{d}t} &\approx \frac{2\mathsf{J}S}{\hbar}(2S_j^{y}-S_{j-1}^{y}-S_{j+1}^{y})\tag{6.15}\\
\frac{\mathrm{d}S_j^{y}}{\mathrm{d}t} &\approx -\frac{2\mathsf{J}S}{\hbar}(2S_j^{x}-S_{j-1}^{x}-S_{j+1}^{x})\tag{6.16}\\
\frac{\mathrm{d}S_j^{z}}{\mathrm{d}t} &\approx 0.\tag{6.17}
\end{align*}
寻找简正模解，令
\begin{equation*}
S_j^{x} = A\,\mathrm{e}^{\mathrm{i}(qja-\omega t)}\tag{6.18}
\end{equation*}
\begin{equation*}
S_j^{y} = B\,\mathrm{e}^{\mathrm{i}(qja-\omega t)}\tag{6.19}
\end{equation*}
q 为波矢。少许代数运算给出 $A=\mathrm{i}B$（$x$ 与 $y$ 方向运动相位差 $\pi/2$），
从而
\begin{equation*}
\hbar\omega = 4\mathsf{J}S(1-\cos qa),\tag{6.20}
\end{equation*}
这就是自旋波的色散关系，绘于图 6.12。自旋波的样子如图 6.13 所示。

\begin{figure}
\centering
\includegraphics[width=0.45\textwidth]{fig6_12.pdf}
\caption{一维自旋链的自旋波色散关系。}
\end{figure}

\begin{figure}
\centering
\includegraphics[width=0.4\textwidth]{fig6_13.pdf}
\caption{自旋链上的自旋波。（a）透视图。（b）俯视图。}
\end{figure}

（2）用量子力学方法导出式 6.20 或许更令人满意，现在重复这一推导。\fn{4}{目前取 $S=\frac{1}{2}$，但很容易推广。}考虑系统基态 $|\Phi\rangle$——所有自旋沿 $+z$：
\begin{center}
$|\Phi\rangle = \cdots\uparrow\uparrow\uparrow\uparrow\uparrow\uparrow\uparrow\uparrow\uparrow\cdots$
\end{center}
一维中海森伯模型的哈密顿量可写为
\begin{equation*}
\hat{\mathcal{H}} = -2\mathsf{J}\sum_i\left[\hat{S}_i^{z}\hat{S}_{i+1}^{z} + \frac{1}{2}(\hat{S}_i^{+}\hat{S}_{i+1}^{-}+\hat{S}_i^{-}\hat{S}_{i+1}^{+})\right],\tag{6.21}
\end{equation*}
于是 $\hat{\mathcal{H}}|\Phi\rangle=-NS^{2}\mathsf{J}|\Phi\rangle$。现在制造一个
激发：翻转格点 j 的自旋，考虑态 $|j\rangle=\hat{S}_j^{-}|\Phi\rangle$——格点 j
自旋翻转了的基态（见图 6.14）：
\begin{center}
$|j\rangle = \cdots\uparrow\uparrow\downarrow\uparrow\uparrow\uparrow\uparrow\uparrow\uparrow\cdots$
\end{center}
翻转一个自旋使系统总自旋改变了 $\frac{1}{2}-(-\frac{1}{2})=1$，因此该激发有整数
自旋、是玻色子。把哈密顿量作用到这个新态上：
\begin{equation*}
\hat{\mathcal{H}}|j\rangle = 2\left[(-NS^{2}\mathsf{J}+2S\mathsf{J})|j\rangle - S\mathsf{J}|j+1\rangle - S\mathsf{J}|j-1\rangle\right],\tag{6.22}
\end{equation*}
它不是 $|j\rangle$ 的常数倍——这个态不是哈密顿量的本征态。然而，可以通过寻找
平面波形式的解来对角化哈密顿量：
\begin{equation*}
|q\rangle = \frac{1}{\sqrt{N}}\sum_j\mathrm{e}^{\mathrm{i}qR_j}|j\rangle.\tag{6.23}
\end{equation*}
态 $|q\rangle$ 本质上是一个离域化（抹开）到所有格点的翻转自旋。由于它由“单个翻转
自旋”的态组成，$|q\rangle$ 自身的总自旋为 $NS-1$。于是容易证明
\begin{equation*}
\hat{\mathcal{H}}|q\rangle = E(q)|q\rangle,\tag{6.24}
\end{equation*}
其中
\begin{equation*}
E(q) = -2NS^{2}\mathsf{J} + 4\mathsf{J}S(1-\cos qa).\tag{6.25}
\end{equation*}
激发的能量为 $\hbar\omega=4JS(1-\cos qa)$，与上面式 6.20 的结果相同。

\begin{figure}
\centering
\includegraphics[width=0.65\textwidth]{fig6_14.pdf}
\caption{态 $|\Phi\rangle$ 由所有沿 $+z$ 方向的自旋组成；态 $|j\rangle$ 是格点 j
处自旋翻转了的基态。}
\end{figure}

\subsection{布洛赫 $T^{3/2}$ 定律}

小 q 时式 6.20 给出
\begin{equation*}
\hbar\omega \approx 2\mathsf{J}Sq^{2}a^{2},\tag{6.26}
\end{equation*}
即 $\omega\propto q^{2}$。三维中态密度为
\begin{equation*}
g(q)\,\mathrm{d}q \propto q^{2}\,\mathrm{d}q\tag{6.27}
\end{equation*}
从而在低温（只有小 q、小 $\omega$ 重要）下
\begin{equation*}
g(\omega)\,\mathrm{d}\omega \propto \omega^{1/2}\,\mathrm{d}\omega\tag{6.28}
\end{equation*}
自旋波的量子化方式与格波相同：后者称声子，前者相应称磁振子。如 6.6.1 节所示，
磁振子自旋为 1、是玻色子。

温度 T 时被激发的磁振子模数 $n_{\mathrm{magnon}}$ 由磁振子态密度对所有频率积分
得到，须乘上玻色因子 $(\exp(\hbar\omega/k_{\mathrm{B}}T)-1)^{-1}$（磁振子是玻色子）：
\begin{equation*}
n_{\mathrm{magnon}} = \int_{0}^{\infty}\frac{g(\omega)\,\mathrm{d}\omega}{\exp(\hbar\omega/k_{\mathrm{B}}T)-1},\tag{6.29}
\end{equation*}
用代换 $x=\hbar\omega/k_{\mathrm{B}}T$ 求值。低温下（三维 $g(\omega)\propto\omega^{1/2}$）
得
\begin{equation*}
n_{\mathrm{magnon}} = \left(\frac{k_{\mathrm{B}}T}{\hbar}\right)^{3/2}\int_{0}^{\infty}\frac{x^{1/2}\,\mathrm{d}x}{\mathrm{e}^{x}-1} \propto T^{3/2}.\tag{6.30}
\end{equation*}
由于每个热激发的磁振子模把总磁化强度减少 S=1，\fn{5}{因为每个磁振子模是一个离域化的单个翻转自旋，见上一节。}低温下自发磁化强度相对 T=0 值的减少为
\begin{equation*}
\frac{M(0)-M(T)}{M(0)} \propto T^{3/2}.\tag{6.31}
\end{equation*}
这一结果称为布洛赫 $T^{3/2}$ 定律（Bloch $T^{3/2}$ law），在低温区与实验数据
符合（见图 6.15）。\fn{6}{布洛赫 $T^{3/2}$ 定律其实只对畴内的自发磁化强度严格成立。图 6.15 的数据实际测量的正是它——它们是用 $\mu$SR 在零外场下测得的局域内场。}磁振子模的能量为
\begin{equation*}
E_{\mathrm{magnon}} = \int_{0}^{\infty}\frac{\hbar\omega\,g(\omega)\,\mathrm{d}\omega}{\exp(\hbar\omega/k_{\mathrm{B}}T)-1} \propto T^{5/2},\tag{6.32}
\end{equation*}
故热容 $C=\partial E_{\mathrm{magnon}}/\partial T$ 也正比于 $T^{3/2}$。

\begin{figure}
\centering
\includegraphics[width=0.5\textwidth]{fig6_15.pdf}
\caption{铁磁体中的自发磁化强度。低温下可用自旋波模型拟合、遵循布洛赫 $T^{3/2}$
律；接近临界温度时磁化强度正比于 $(T-T_{\mathrm{C}})^{\beta}$（$\beta$ 为临界
指数）。两种行为都不能在整个温度范围拟合真实数据。数据取自一种有机铁磁体，
$T_{\mathrm{C}}\approx0.67$ K，$\beta\approx0.36$——适用于三维海森伯模型。}
\end{figure}

\subsection{Mermin--Wagner--Berezinskii 定理}

然而在一维与二维，式 6.29 的积分发散，故对各向同性的一维、二维海森伯模型，
一切 $T>0$ 都有 $M\to0$。\fn{7}{对不同的维数，式 6.29 中的函数 $g(\omega)$ 对 $\omega$ 有不同的依赖。}有限温度下产生的自旋波数目发散，因此自发铁磁性不可能。这一结果由
Mermin 与 Wagner 于 1966 年首先证明，Berezinskii 独立证明。二维体系（具有连续
对称性）中长程序的缺失常称为 Mermin--Wagner--Berezinskii 定理。

该结果只适用于各向同性的海森伯铁磁体。它具有旋转对称——所有自旋方向可以整体
旋转而不费任何额外能量。这意味着长波长激发——自旋态在相当长距离上偏离基态值
——几乎不费能量。\fn{8}{若自旋波波长足够长，交换能量代价 $\propto\int(|\nabla M_x|^{2}+|\nabla M_y|^{2})\,\mathrm{d}x\,\mathrm{d}y$（二维）被最小化。见习题 6.7。}于是自旋涨落以极小的能量代价即可激发；在一维与二维，它们摧毁长程序。然而若存在
显著各向异性，把自旋从基态值转开将有能量代价。事实证明，容忍这些涨落所要付出的
各向异性能量代价随激发半径 R 的平方增长，故各向异性能量会抑制除最小者之外的所有
非线性涨落。\fn{9}{这一想法在习题 6.7 中探讨。}正是这种对称性破缺场的存在，可以稳定二维体系中的长程序。
真实体系中自旋间还存在偶极相互作用——虽远弱于交换相互作用，却是各向异性的，
能以类似方式抑制涨落的生长。

实验上观察到，一些超薄磁性薄膜可以表现自发铁磁性。有时各向异性使自旋垂直于膜面
在能量上有利（系统被称为具有垂直易磁化轴）。此时各向异性在自旋波激发谱中（基态
与激发态之间）打开能隙，长程磁有序可以在有限温度存在。偶极相互作用或膜面内自旋
能量的各向异性也能导致类似效果、稳定磁化强度。

\subsection{自旋波的测量}

用于磁结构测定的弹性中子散射已在 5.7.2 节讨论。自旋波色散可以用非弹性中子散射
测量：此时入射中子波矢的大小不再等于散射中子波矢 $k'$ 的大小，中子能量也从
\begin{equation*}
E = \frac{\hbar^{2}k^{2}}{2m_{\mathrm{n}}}\tag{6.33}
\end{equation*}
变为
\begin{equation*}
E' = \frac{\hbar^{2}k'^{2}}{2m_{\mathrm{n}}}\tag{6.34}
\end{equation*}
因为中子在样品中产生了一个能量 $\hbar\omega$、波矢 q 的激发。能量与动量守恒给出
\begin{equation*}
E = E' + \hbar\omega\tag{6.35}
\end{equation*}
\begin{equation*}
\mathbf{k} = \mathbf{k}' + \mathbf{q} + \mathbf{G},\tag{6.36}
\end{equation*}
G 是倒格矢波矢，\fn{10}{加入 G 是必要的：磁振子的色散关系与声子一样，在倒格子上是周期的。}于是测得 $\mathbf{k}$、$\mathbf{k}'$、E、$E'$ 便能确定 $\omega$ 与 q。
中子的能量与原子过程及电子过程的能量相近，即 meV 到 eV 范围。

因此可以探测从量子隧穿的 $\mu$eV，经过分子的平移、转动、振动与晶格模式，直到
材料电子结构中 eV 跃迁的能量尺度；磁振子能量典型在 $10^{-3}$--$10^{-2}$ eV，
因此可以有效地用非弹性中子散射测量。

实验通常用中子三轴谱仪（triple-axis spectrometer，见图 6.16）进行——如此命名
是因为单色器、样品与分析器晶体的角度都可以独立变化。这允许测量对应大范围可能
$\omega$ 与 $\mathbf{q}$ 的散射中子，特别地可以作固定 $\omega$ 变 q（或反之）的
扫描。用该技术获得的数据示例：图 6.17 为 $\mathrm{Co}_{0.92}\mathrm{Fe}_{0.08}$，
图 6.18 为铁磁氧化物 $\mathrm{La}_{0.7}\mathrm{Pb}_{0.3}\mathrm{MnO}_3$
（它表现庞磁电阻，见 8.9.5 节）。

\begin{figure}
\centering
\includegraphics[width=0.45\textwidth]{fig6_16.pdf}
\caption{中子三轴谱仪示意。设置旋转屏蔽与单色器晶体的角度可选定入射中子能量；
样品也可旋转；分析器晶体也可旋转，从而测量散射中子的能量。}
\end{figure}

\begin{figure}
\centering
\includegraphics[width=0.45\textwidth]{fig6_17.pdf}
\caption{$\mathrm{Co}_{0.92}\mathrm{Fe}_{0.08}$ 合金中自旋波的磁振子能谱，
室温测得（Sinclair 与 Brockhouse 1960）。图中 k 是正文中的磁振子波矢 q。}
\end{figure}

\begin{figure}
\centering
\includegraphics[width=0.45\textwidth]{fig6_18.pdf}
\caption{铁磁氧化物 $\mathrm{La}_{0.7}\mathrm{Pb}_{0.3}\mathrm{MnO}_3$ 中
10 K 非弹性中子散射测得的自旋波色散；横轴为磁振子波矢。取自 Perring 等 1996。}
\end{figure}

\section{磁畴}

若宏观体系的不同区域以不同方式打破对称，则这些区域的界面上刚性可能失效。一般
预期存在畴壁、缺陷、涡旋、位错与其他奇点。铁磁体中最重要的奇点是畴壁（domain
wall）。

外斯首先提出铁磁体包含许多称为磁畴（domains）的小区域，每个畴内局域磁化达到
饱和值；不同畴的磁化方向不必平行；畴由畴壁分隔。磁畴的存在解释了令人惊讶的
观察：某些（磁“软”的\fn{11}{即磁‘软’的材料，见第 6.7.9 节。}）铁磁样品在室温下仅用极弱的磁场（低至 $10^{-6}$ T）就能使整块样品达到饱和磁化强度（相应
$\mu_0M\sim1$ T）。如此低的外场对顺磁体几乎没有效果。\fn{12}{若某顺磁体的 $\chi\sim10^{-3}$（见表 2.1），则 $10^{-6}$ T 的外场只能产生 $\mu_0M\sim10^{-9}$ T 的磁化强度。}铁磁样品中的巨大效应在于：外场无须宏观地把磁矩排序（它们早已有序！），
只须使各磁畴排齐——这可以通过能量上无痛的畴壁移动过程实现。\fn{13}{我们将看到，某些铁磁材料中畴壁的移动实际上是耗能的；这类材料需要强外场才能从一个磁态切换到另一个。}

同一铁磁样品中，零外场下磁化强度也可能为零——这也是磁畴的表现：每个畴内磁化
饱和，但各畴磁化方向安排得使样品净磁化强度为零。

\subsection{畴壁}

相邻畴之间的边界称为畴壁。畴壁可按两畴磁化强度间的夹角分类（见图 6.19）：180°
畴壁分隔磁化相反的畴；90° 畴壁分隔磁化垂直的畴。

最常见的 180° 壁是布洛赫壁（Bloch wall，见图 6.20(a)）：磁化在平行于壁面的平面内
旋转。另一种可能组态是奈尔壁（Néel wall，见图 6.20(b)）：磁化在垂直于壁面的平面内
旋转。下面尝试计算布洛赫壁的畴壁宽度。

铁磁体中转动相邻自旋要费能量。夹角为 $\theta$ 的两个自旋 $S_1$、$S_2$
（图 6.21）能量为 $-2\mathsf{J}S_1\cdot S_2=-2\mathsf{J}S^{2}\cos\theta$；
$\theta=0$ 时能量 $-2\mathsf{J}S^{2}$。于是 $\theta\neq0$ 的能量代价在 $\theta\ll1$
时近似为\fn{14}{利用 $\theta\ll1$ 时 $\cos\theta\approx1-\theta^{2}/2$。}$\mathsf{J}S^{2}\theta^{2}$。布洛赫壁中自旋跨 N 个格点转过 $\pi$ 角（图 6.20(a)），一条自旋线的能量代价是 N 个
$\mathsf{J}S^{2}\theta^{2}$ 贡献之和，$\theta=\pi/N$，即 $\mathsf{J}S^{2}\pi^{2}/N$。
布洛赫壁中是自旋平面（见图 6.20），因此我们关心的是布洛赫壁单位面积的能量
$\sigma_{\mathrm{BW}}$。一平方米的壁上有 $1/a^{2}$ 条我们计算过的自旋线。故
\begin{equation*}
\sigma_{\mathrm{BW}} = \mathsf{J}S^{2}\frac{\pi^{2}}{Na^{2}}\tag{6.37}
\end{equation*}
它在 $N\to\infty$ 时趋于零。这一结果似乎表明：一旦形成畴壁，它就会自己解开、在
整个系统中不断长大——因为自旋相互扭曲要费能量，除非有别的相互作用阻止，它们都
会解开。这个别的相互作用就是磁晶各向异性（magnetocrystalline anisotropy）。

\begin{figure}
\centering
\includegraphics[width=0.4\textwidth]{fig6_19.pdf}
\caption{（a）180° 畴壁。（b）90° 畴壁。}
\end{figure}

\begin{figure}
\centering
\includegraphics[width=0.45\textwidth]{fig6_20.pdf}
\caption{（a）布洛赫壁，（b）奈尔壁。}
\end{figure}

\begin{figure}
\centering
\includegraphics[width=0.4\textwidth]{fig6_21.pdf}
\caption{彼此成 $\theta$ 角的两个自旋，产生能量代价 $\frac{1}{2}\mathsf{J}S^{2}\theta^{2}$。}
\end{figure}

\subsection{磁晶各向异性}

晶体具有易磁化轴与难磁化轴：沿某些晶体学方向容易磁化晶体，沿另一些则较难
（Fe、Co、Ni 单晶的情形见图 6.22）。例如在 Co 中，这一各向异性导致形如
\begin{equation*}
E = K_1\sin^{2}\theta + K_2\sin^{4}\theta\tag{6.38}
\end{equation*}
的附加能量，$K_1$、$K_2$ 是各向异性常数，$\theta$ 是磁化强度与六角密堆面堆叠
方向之间的夹角。由于这些常数为正，磁化沿堆叠方向时能量最低。式 6.38 中 E、$K_1$、
$K_2$ 是能量密度（以 J\,m$^{-3}$ 量度）。各向异性常数强烈依赖温度。

\begin{figure}
\centering
\includegraphics[width=0.55\textwidth]{fig6_22.pdf}
\caption{Fe、Co、Ni 沿不同方向加场时的磁化强度，显示各向异性。取自 Honda 与 Kaya
1926、Kaya 1928。}
\end{figure}

式 6.38 适用于单轴各向异性——能量只依赖与单一轴的夹角（Co 中即六角密堆面的堆叠
轴）。立方体系中恰当的表达是
\begin{equation*}
E = K_1(m_x^{2}m_y^{2}+m_y^{2}m_z^{2}+m_z^{2}m_x^{2}) + K_2m_x^{2}m_y^{2}m_z^{2}+\cdots,\tag{6.39}
\end{equation*}
$\mathbf{m}=(m_x,m_y,m_z)=\mathbf{M}/|\mathbf{M}|$。球坐标下为
\begin{equation*}
E = K_1\left(\frac{1}{4}\sin^{2}\theta\sin^{2}2\phi+\cos^{2}\theta\right)\sin^{2}\theta + \frac{K_2}{16}\sin^{2}2\phi\sin^{2}2\theta\sin^{2}\theta+\cdots.\tag{6.40}
\end{equation*}
各向异性能量源于自旋--轨道相互作用与角动量的部分淬灭。各向异性能量通常在
$10^{2}$--$10^{7}$ J\,m$^{-3}$，相当于每原子 $10^{-8}$--$10^{-3}$ eV。对称性低的
（磁性离子）晶格各向异性能量较大，高对称者较小。例如立方 Fe、Ni 的 $K_1$ 分别为
$4.8\times10^{4}$ 与 $-5.7\times10^{3}$ J\,m$^{-3}$，六角 Co 的 $K_1=5\times10^{5}$
J\,m$^{-3}$；低对称的永磁材料 $\mathrm{Nd_2Fe_{14}B}$ 与 $\mathrm{SmCo_5}$ 的
$K_1$ 分别为 $5\times10^{6}$ 与 $1.7\times10^{7}$ J\,m$^{-3}$。

另一项附加能量来自与样品形状相联系的退磁能，称为形状各向异性（shape anisotropy）。
薄膜中形状各向异性 $\frac{1}{2}\mu_0M^{2}\cos^{2}\theta$（$\theta$ 为膜法线与 M
的夹角）使磁化强度保持在膜面内以节省能量。

\subsection{畴壁宽度}

铁磁体的磁畴中，磁化偏好沿易方向排列；但在畴之间的畴壁里它必须转动，一部分沿
难轴、耗费能量。设各向异性能量密度取简单形式 $E=K\sin^{2}\theta$（K 为各向异性
常数），便可轻松求出各向异性能量对布洛赫壁的贡献。取 K>0，自旋偏好沿 $\theta=0$
或 $\theta=\pi$。把 N 个自旋的贡献相加，并把求和换成积分（连续极限）以简化：
能量密度贡献为
\begin{equation*}
\sum_{i=1}^{N}K\sin^{2}\theta_i \approx \frac{N}{\pi}\int_{0}^{\pi}K\sin^{2}\theta\,\mathrm{d}\theta = \frac{NK}{2}.\tag{6.41}
\end{equation*}
壁单位面积的能量贡献于是为 NKa/2。畴壁单位面积的总能量（含交换能贡献（式 6.37）
与各向异性能量贡献（式 6.41））为
\begin{equation*}
\sigma_{\mathrm{BW}} = \mathsf{J}S^{2}\frac{\pi^{2}}{Na^{2}} + \frac{NKa}{2}.\tag{6.42}
\end{equation*}
这给出所需的行为：一项正比于 1/N（趋向解开壁、使之变大），另一项正比于 N（趋向
收紧壁、使之变小）。用 $\mathrm{d}E_{\mathrm{BW}}/\mathrm{d}N=0$ 求平衡组态，得
\begin{equation*}
N = \pi S\sqrt{2\mathsf{J}/Ka^{3}}\tag{6.43}
\end{equation*}
布洛赫壁宽度为
\begin{equation*}
\delta = Na = \pi S\sqrt{\frac{2\mathsf{J}}{Ka}},\tag{6.44}
\end{equation*}
a 是晶格间距。$\mathsf{J}$ 越大壁越厚，K 越大壁越薄。畴壁单位面积的能量为
\begin{equation*}
\sigma_{\mathrm{BW}} = \pi S\sqrt{\frac{2\mathsf{J}K}{a}}.\tag{6.45}
\end{equation*}
用式 4.19 定义的参数 A（立方晶体）表示，畴壁厚度为
\begin{equation*}
\delta = \pi\sqrt{\frac{A}{K}},\tag{6.46}
\end{equation*}
单位面积能量为
\begin{equation*}
\sigma_{\mathrm{BW}} = \pi\sqrt{AK}.\tag{6.47}
\end{equation*}

\subsection{磁畴的形成}

既然形成畴壁要付出能量，人们或许奇怪：为什么一开始会形成不止一个畴？原因是磁畴
的形成节省与偶极场相联系的能量。由于 $\nabla\cdot\mathbf{H}=-\nabla\cdot\mathbf{M}$
（见附录 B），凡 M 中断或开始之处（例如样品边缘），磁场的散度不为零，产生退磁场，
充满空间，每立方米耗费 $B^{2}/2\mu_0$ 焦耳的能量。与退磁场相联系的能量称为退磁能、
静磁能或偶极能：
\begin{equation*}
-\frac{\mu_0}{2}\int_{V}\mathbf{M}\cdot\mathbf{H}_{\mathrm{d}}\,\mathrm{d}\tau.\tag{6.48}
\end{equation*}
$\mathbf{H}_{\mathrm{d}}$ 是退磁场，积分遍及样品体积（见附录 D）。对沿主轴磁化的
椭球样品，此能量化为
\begin{equation*}
\frac{\mu_0}{2}NM^{2}V\tag{6.49}
\end{equation*}
N 是退磁因子，V 是样品体积。

把样品分成磁畴可以节省这一偶极能，但每创造一个畴都因畴壁而耗能。因此，正如畴壁
尺寸是交换能与各向异性能量的平衡，磁畴的形成是退磁场代价与畴壁代价的平衡。

\begin{figure}
\centering
\includegraphics[width=0.3\textwidth]{fig6_23.pdf}
\caption{铁磁样品的三种畴结构选择：（a）均匀磁化；（b）分成两个畴；（c）简单
封闭畴结构。}
\end{figure}

图 6.23 给出铁磁样品磁畴结构的三种选择：图 6.23(a) 的单畴结构没有畴壁但偶极能
巨大；把样品分成两畴（图 6.23(b)）可以降低偶极能，代价是引入畴壁；图 6.23(c) 的
封闭畴（closure domain）结构消除偶极能但引入若干畴壁。

偶极能还能决定可形成的畴壁类型。布洛赫壁在块体中更受青睐，因为其偶极能更小；
奈尔壁则常见于薄膜（把自旋转出膜面要付偶极能代价）。

\subsection{磁化过程}

图 6.24 给出测量铁磁样品磁化强度随外加磁场变化的典型磁滞回线。若样品被外场磁化
到饱和磁化强度 $M_{\mathrm{s}}$，外场降为零时磁化强度降为剩余磁化强度
（remanent magnetization）$M_{\mathrm{r}}$；需要等于矫顽场（coercive field）
$H_{\mathrm{c}}$ 的磁场才能把磁化强度切换到相反方向。参数 $M_{\mathrm{r}}$ 与
$H_{\mathrm{c}}$ 可用来表征铁磁体。

\begin{figure}
\centering
\includegraphics[width=0.42\textwidth]{fig6_24.pdf}
\caption{磁滞回线，标出饱和磁化强度 $M_{\mathrm{s}}$、剩余磁化强度 $M_{\mathrm{r}}$
与矫顽场 $H_{\mathrm{c}}$。}
\end{figure}

\begin{figure}
\centering
\includegraphics[width=0.5\textwidth]{fig6_25.pdf}
\caption{外场对铁晶须表面磁畴图样的影响：随外场 B 从 0 增大到最大值，畴壁发生
移动。}
\end{figure}
\bio{John Kerr\\（1824--1907）\quad Heinrich Barkhausen\\（1881--1956）}

铁磁体磁化时长度会略微改变——这一效应称为磁致伸缩（magnetostriction），表明
材料的磁学与力学行为相互联系（该联系称为磁弹耦合）。晶体形变是因为可以降低其
各向异性能量——后者可依赖晶体应变。例如立方体系中，若轻微偏离严格立方对称所节省
的各向异性能量超过弹性能量的代价，则形变在能量上有利。

退磁铁磁体被磁化时发生多种过程。首先是畴壁移动：相对外场取向有利的畴吞并不利的
畴（图 6.25 示意这一过程）。更高场下可发生畴转动：各向异性能量被压倒，磁化强度
突然从原磁化方向转到最靠近场方向的晶体学易轴。最高磁场下的最后一个畴过程是各畴
的一致转动——无论易轴难轴，转到与磁场排齐。

畴壁在磁性材料中的运动细节依赖于材料的冶金学性质。由于磁弹耦合，畴壁可被材料中
的应变、表面与杂质钉扎（pinning）；畴壁钉扎因此提高矫顽力。铁磁体的磁化强度也因
畴界运动而以一系列不连续台阶变化，磁化曲线上有时能看到极小的台阶——这称为巴克
豪森效应（Barkhausen effect）。磁化过程中这些突然的不连续变化及其与系统弹性模式
的耦合，有时伴随低水平的声发射，称为磁声发射（magnetoacoustic emission）。

\subsection{磁畴的观察}

观察磁畴有多种技术。比脱粉纹法（Bitter powder technique）把细磁性颗粒的胶体
悬浮液滴在样品表面；磁性颗粒聚集在畴壁附近——那里有吸引它们的强局域场——可用
光学显微镜观察。也可以用偏振光从抛光磁化样品表面的反射观察磁畴，称为克尔显微镜
（Kerr microscopy）：入射光被偏振，其偏振面被磁化样品的反射所转动（即克尔效应，
见 8.8 节）；转角可用第二个偏振器（检偏器）测量。若把光束聚焦成一点并扫过磁性
样品表面，则在检偏器的特定设置下，一些畴显得亮、另一些暗。若样品透明，实验可在
透射下进行（法拉第效应）。

电子显微术同样有效：电子束被内场的洛伦兹力偏转，畴之间的边界（扫描经过时偏转
可以突然变号）显示得特别清楚。这称为洛伦兹显微术（Lorentz microscopy）。

磁力显微术（magnetic force microscopy\fn{15}{磁力显微镜以著名的原子力显微镜为基础。}）用一个带磁化针尖的微悬臂扫过样品表面：畴壁产生的杂散场使悬臂弯曲，
这一弯曲（例如通过悬臂共振频率的变化探测）用于成像表面。

\subsection{小磁性颗粒}

许多铁磁样品在零外场下能量最低的态是退磁态，总磁矩为零。若减小样品尺寸，表面能
（如畴壁能）相对体积能（如退磁能）变得越发昂贵。\fn{16}{这是因为表面能按（样品尺寸）$^2$ 标度，而体积能按（样品尺寸）$^3$ 标度。}于是会到达某个临界尺寸：低于它，去掉畴壁、让样品成为单个磁化
畴在能量上有利——它因此像一块小永磁体。

下面的例子估计单畴颗粒的临界尺寸。计算将假定该临界尺寸仍大于畴壁宽度。

\begin{example}{6.1}
考虑半径 r 的铁磁颗粒。单畴态（图 6.26(a)）的能量用式 D.12 并取球的
$N=\frac{1}{3}$：\mnote{球的体积 $V=\frac{4}{3}\pi r^{3}$。}
\begin{equation*}
E_{(a)} = \frac{1}{6}\mu_0M^{2}V = \frac{2}{9}\mu_0\pi M^{2}r^{3}.\tag{6.50}
\end{equation*}
图 6.26(b) 的态（适用于立方各向异性）近似消除全部退磁能，但引入面积为
$2\times\pi r^{2}$ 的 90° 畴壁，能量代价
\begin{equation*}
E_{(b)} = 2\pi r^{2}\sigma_{\mathrm{W}}^{90^{\circ}},\tag{6.51}
\end{equation*}
$\sigma_{\mathrm{W}}^{90^{\circ}}$ 是 90° 畴壁单位面积的能量代价。当
$E_{(a)}<E_{(b)}$，即
\begin{equation*}
r < \frac{9\sigma_{\mathrm{W}}^{90^{\circ}}}{\mu_0M^{2}}.\tag{6.52}
\end{equation*}
时，图 6.26(a) 的态比图 6.26(b) 更有利。

图 6.26(c) 的态（适用于单轴各向异性）较难考虑，因为两个畴都不是椭球形、仍有
一些退磁能。不过粗略地说，该组态的退磁能是图 6.26(a) 的一半（畴只有一半大小）；
畴壁代价为 $\pi r^{2}\sigma_{\mathrm{W}}^{180^{\circ}}$（$\sigma_{\mathrm{W}}^{180^{\circ}}$
是 180° 畴壁单位面积的能量）。于是
\begin{equation*}
E_{(c)} = \frac{1}{9}\mu_0\pi M^{2}r^{3} + \pi r^{2}\sigma_{\mathrm{W}}^{180^{\circ}}\tag{6.53}
\end{equation*}
故当 $E_{(a)}<E_{(c)}$，即
\begin{equation*}
r < \frac{9\sigma_{\mathrm{W}}^{180^{\circ}}}{\mu_0M^{2}}.\tag{6.54}
\end{equation*}
时，图 6.26(a) 的态比图 6.26(c) 更有利。

若 $\sigma_{\mathrm{W}}^{180^{\circ}}\sim10^{-2}$ J\,m$^{-2}$、$\mu_0M\sim1$ T，
则临界半径 $\sim10^{-7}$ m。注意：若该值小于畴壁宽度，计算便失去意义。
\end{example}

\begin{figure}
\centering
\includegraphics[width=0.3\textwidth]{fig6_26.pdf}
\caption{小球形铁磁颗粒的三种可能磁化组态。}
\end{figure}

\subsection{Stoner--Wohlfarth 模型}

可以通过 Stoner--Wohlfarth 模型计算单畴颗粒的磁化曲线。因为颗粒是单个磁畴，
没有畴壁，只须考虑一致的畴转动。

考虑磁场 H 中的单畴磁性颗粒，H 与其单轴各向异性易轴成 $\theta$ 角。若颗粒磁化
强度与磁场方向成 $\phi$ 角，系统的能量密度为
\begin{equation*}
E = K\sin^{2}(\theta-\phi) - \mu_0HM_{\mathrm{s}}\cos\phi.\tag{6.55}
\end{equation*}
对任意给定的外加磁场，能量取极小即得磁化强度的方向。

图 6.27 给出结果：外加场用无量纲参数
\begin{equation*}
h = \frac{\mu_0M_{\mathrm{s}}H}{2K}\tag{6.56}
\end{equation*}
表示，对两个外场方向绘出。该模型对 $\theta=0$ 与 $\theta=\pi/2$ 有解析解，但图中
展示的是如何数值求解。每种情形都画出在能量面（能量作为 $\phi$ 与 h 的函数）上降低
能量的磁化方向轨迹，并给出相应的磁滞回线。图 6.28 给出更大范围 $\theta$ 值的
磁滞回线，以及多晶方向平均的计算结果（适用于单畴颗粒阵列，忽略颗粒间相互作用）。
该模型表明：即使在没有畴壁钉扎这类不可逆效应的体系中，系统中的各向异性也能导致
磁滞。

\begin{figure}
\centering
\includegraphics[width=0.4\textwidth]{fig6_27.pdf}
\caption{Stoner--Wohlfarth 模型中（a）$\theta=90^{\circ}$ 与（c）$\theta=30^{\circ}$
的磁滞回线。它们可以由在能量面（作为 $h$ 与 $\phi$ 的函数，（b）$\theta=90^{\circ}$
与（d）$\theta=30^{\circ}$）上寻找能量极小点得到。}
\end{figure}

\begin{figure}
\centering
\includegraphics[width=0.4\textwidth]{fig6_28.pdf}
\caption{Stoner--Wohlfarth 模型的磁滞回线：（a）$\theta=0^{\circ}$（粗线）、
$5^{\circ}$、$15^{\circ}$、$30^{\circ}$；（b）$\theta=45^{\circ}$、$60^{\circ}$、
$75^{\circ}$ 与 $90^{\circ}$（粗线）。（c）多晶平均的计算磁滞回线。}
\end{figure}

\subsection{软磁与硬磁材料}

铁磁体沿磁滞回线绕行一周所耗散（为热）的能量正比于回线面积。面积小者称磁软
（magnetically soft）材料；面积大者称磁硬（magnetically hard）材料。磁场循环时
畴壁穿越样品；可以按照畴在样品中移动的难易区分这两大类铁磁材料。

（1）软磁材料容易磁化。软磁材料用于变压器线圈、发电机与电动机：磁化强度须每秒
反转许多次，每周期耗散的能量须最小。软材料有宽畴壁（各向异性能量 K 小）因而容易
移动，矫顽场小。低磁致伸缩也常是想要的——使内应变不诱生局域各向异性能量。例子是
坡莫合金（permalloy，一种商用 Ni/Fe 合金，另加成分），矫顽场
$B_{\mathrm{c}}\sim2\times10^{-7}$ T。

（2）硬磁材料难以磁化，因而难以退磁。硬磁体用作永磁体（如扬声器背部、电动机，
当然还有你冰箱的门上！）以及磁带记录（粉体形式）。这些应用中磁化强度须尽可能
长久保持：磁滞回线每周期的能量耗散越大越好，磁化便不会自发消失。硬磁体磁滞大、
畴壁窄（K 大），畴壁钉扎容易实现。大的离子磁矩与大的晶体场有利于硬磁性质，合适
的材料常涉及稀土，如 $\mathrm{Nd_2Fe_{14}B}$：$T_{\mathrm{C}}=585$ K，矫顽场
$B_{\mathrm{c}}=1.2$ T。

硬磁材料的一个重要应用是磁记录。磁性材料可记录的比特面密度近年来猛增，反映了
对磁性材料的理解与制备的稳步进步。大致数字：1960 年 $\sim0.1$ Mbits\,in$^{-2}$，
1980 年 $\sim10$ Mbits\,in$^{-2}$，2000 年 $\sim25$ Gbits\,in$^{-2}$（1 bit\,in$^{-2}$
即每平方英寸存储一个比特的‘0’或‘1’——这个最现代的产业却还没把单位现代化！）
随着未来面密度继续增大，磁矩的超顺磁涨落（即磁性颗粒的热涨落，见 8.3 节）将变得
重要。

传统上记录磁带与磁盘使用 $\mathrm{Fe_2O_3}$ 小颗粒。高保真录音与录像带性能的
改进基于高矫顽力颗粒：掺 Co 的 $\mathrm{Fe_2O_3}$、$\mathrm{CrO_2}$、金属颗粒
（通常为 Fe）以及钡铁氧体（BaO·6Fe$_2$O$_3$，见 5.3 节）。硬盘通常是一大片刚性
盘基，其上用真空沉积或溅射镀薄膜，通常是 Co/Cr 合金。硬盘高速旋转，读头横越盘面
（可以悬浮在盘面上方一微米的几分之一处），能访问盘面的任何部位。

\section*{延伸阅读}

\begin{itemize}[itemsep=0.2em]
\item 对称性破缺概念的更详细描述见 P.~W.~Anderson,《Basic notions of condensed
matter physics》, Addison-Wesley 1984。
\item 相变的统计力学见 J.~Yeomans,《Statistical mechanics of phase transitions》,
OUP 1992。
\item 关于临界现象也非常有帮助的是 M.~F.~Collins,《Magnetic critical scattering》,
OUP 1989。
\item 对相变与对称性破缺现象（尤其侧重软凝聚态物质）非常透彻的描述是 P.~M.~Chaikin
与 T.~C.~Lubensky,《Principles of condensed matter physics》, CUP 1995。
\item J.~J.~Binney, N.~J.~Dowrick, A.~J.~Fisher 与 M.~E.~J.~Newman,《The theory
of critical phenomena》, OUP 1992，包含对重正化群方法的精彩论述。
\item 关于磁畴与磁性的信息可见 G.~Bertotti,《Hysteresis in magnetism》, Academic
Press 1998 与 A.~Hubert 及 R.~Sch\"{a}fer,《Magnetic domains》, Springer 1998。
\end{itemize}

\section*{本章参考文献}

\begin{itemize}[itemsep=0.2em]
\item K.~Honda 与 S.~Kaya, Sci.~Rep.~Tohoku Univ., \textbf{15}, 721 (1926)；
S.~Kaya, Sci.~Rep.~Tohoku Univ., \textbf{17}, 639, 1157 (1928)。
\item 有用的一般参考书：D.~Craik,《Magnetism: principles and applications》, Wiley
1995；J.~Crangle,《Solid state magnetism》, Edward Arnold 1991。
\item T.~G.~Perring, G.~Aeppli, S.~M.~Hayden, S.~A.~Carter, J.~P.~Remeika 与
S-W.~Cheong, Phys.~Rev.~Lett., \textbf{77}, 711 (1996)。
\item R.~N.~Sinclair 与 B.~N.~Brockhouse, Phys.~Rev., \textbf{120}, 1638 (1960)。
\end{itemize}

\section*{习题}

\exer{6.1} N 个自旋的一维铁磁链由伊辛哈密顿量
\begin{equation*}
\hat{\mathcal{H}} = -2\mathsf{J}\sum_{i=1}^{N-1}\hat{S}_i^{z}\hat{S}_{i+1}^{z} - 2\mathsf{J}\hat{S}_N^{z}\hat{S}_1^{z}.
\end{equation*}
描述（最后一项给出周期性边界条件）。

证明：引入新算符
\begin{equation*}
\zeta_i = 4S_i^{z}S_{i+1}^{z}
\end{equation*}
（本征值 +1 或 $-1$）可求得该系统的配分函数，且配分函数为
$Z=(2\cosh(\mathsf{J}/2k_{\mathrm{B}}T))^{N}$。

用这些结果求 $N\to\infty$ 时链的每自旋热容表达式。推出热容的低、高温行为并画出
热容随温度变化的草图。讨论结果。

\exer{6.2} 单轴铁磁体由哈密顿量
\begin{equation*}
\hat{\mathcal{H}} = -\sum_{ij}\mathsf{J}_{ij}\hat{\mathbf{S}}_i\cdot\hat{\mathbf{S}}_j - \sum_{ij}K_{ij}\hat{S}_i^{z}\hat{S}_j^{z}.\tag{6.57}
\end{equation*}
描述。（a）证明所有自旋沿 z 轴完全排齐的态是哈密顿量的本征态。

（b）求作为波矢 q 的函数的自旋波谱表达式。

（c）把表达式简化到 $\mathsf{J}_{ij}$、$K_{ij}$ 只限最近邻（$\mathsf{J}_0$、
$K_0$）的情形，铁磁体分别为（i）一维链；（ii）二维正方晶格；（iii）三维体心立方
材料。

\exer{6.3} 用习题 6.2 在小波矢下的结果，推出海森伯模型（$K_0=0$）与伊辛模型
（$\mathsf{J}_0=0$）在习题 6.2 的结构（i）、（ii）、（iii）下自旋波数目随温度的
低温依赖。证明你关于情形（i）伊辛模型的结果与习题 6.1 一致，并证明一、二维海森伯
模型在绝对零度以上没有长程磁有序。

\exer{6.4} 把自旋波理论用于反铁磁体，证明色散关系为
\begin{equation*}
\hbar\omega = 4|\mathsf{J}S\sin(qa)|.\tag{6.58}
\end{equation*}
进而证明：反铁磁体在小 q 时 $\omega\propto|q|$，而铁磁体 $\omega\propto q^{2}$。

\exer{6.5} 磁场 H 中铁磁体的朗道理论给出自由能
\begin{equation*}
F(M) = F_0 + a_0(T-T_{\mathrm{C}})M^{2} + bM^{4} - \mu_0MH\tag{6.59}
\end{equation*}
$a_0$ 与 b 为正常数。证明这意味着
\begin{equation*}
M^{2} = u + v\frac{H}{M}\tag{6.60}
\end{equation*}
u 与 v 是应由你确定的常数。对恰好高于 $T_{\mathrm{C}}$、恰好低于 $T_{\mathrm{C}}$
与恰在 $T_{\mathrm{C}}$ 的温度，把 $M^{2}$ 对 H/M 作草图，说明该方法如何用于确定
$T_{\mathrm{C}}$。这就是阿罗特图（Arrott plot，$M^{2}$ 对 H/M 的图）的原理。

\exer{6.6} 求 Stoner--Wohlfarth 模型在 $\theta=0$ 与 $\theta=\pi/2$ 的解析解。

\exer{6.7} 为阐明 Mermin--Wagner--Berezinskii 定理背后的思想，考虑二维中的一个
激发（见图 6.29）：磁化强度作为位置 $\mathbf{r}=(x,y)$ 的函数大小恒定 $|\mathbf{M}|$
而方向变化，由矢量 $\mathbf{M}(\mathbf{r})$ 给定：
\begin{equation*}
\mathbf{M}(\mathbf{r}) = \begin{cases} |\mathbf{M}|\left(\hat{\mathbf{y}}\cos\dfrac{\pi r}{R} - \hat{\mathbf{x}}\sin\dfrac{\pi r}{R}\right) & r<R\\[8pt] -|\mathbf{M}|\hat{\mathbf{y}} & r>R \end{cases}\tag{6.61}
\end{equation*}

\begin{figure}
\centering
\includegraphics[width=0.45\textwidth]{fig6_29.pdf}
\caption{二维磁体中的一个激发。}
\end{figure}

解释为什么该激发的交换能量代价正比于
\begin{equation*}
\int\left[(\nabla M_x)^{2}+(\nabla M_y)^{2}\right]\mathrm{d}x\,\mathrm{d}y.\tag{6.62}
\end{equation*}
并证明它是常数（$\pi^{3}|\mathbf{M}|^{2}$）、与 R 无关。这意味着大尺度涨落与
小尺度涨落相比没有额外能量代价——从而摧毁一切长程序。本题基于 A.~S.~Arrott 与
B.~Heinrich 的文章中的例子，Journal of Magnetism and Magnetic Materials
\textbf{93}, 571 (1991)。

\exer{6.8} 数据以 25 Gbits\,in$^{-2}$ 的面密度写到磁记录盘上。用一些有根据的估计，
把这个面密度换算成其他单位：（a）每个氢原子截面积上的比特数；（b）每张邮票可存
《莎士比亚全集》（按其信息量计）的份数。把这一信息存储密度与人类基因组比较。

\exer{6.9} 考虑磁场中各向同性立方铁磁体（晶格参数 a），哈密顿量
\begin{equation*}
\hat{\mathcal{H}} = -\mathsf{J}\sum_{\langle ij\rangle}\hat{\mathbf{S}}_i\cdot\hat{\mathbf{S}}_j + g\mu_{\mathrm{B}}\sum_j\mathbf{B}\cdot\hat{\mathbf{S}}_j.\tag{6.63}
\end{equation*}
证明 $qa\ll1$ 时其自旋波色散为
\begin{equation*}
\hbar\omega(q) = g\mu_{\mathrm{B}}B + 2\mathsf{J}Sq^{2}a^{2}\tag{6.64}
\end{equation*}

\exer{6.10} 证明哈密顿量
\begin{equation*}
\hat{\mathcal{H}} = -\sum_{\mathbf{n},\delta}\mathsf{J}(\delta)\hat{\mathbf{S}}_{\mathbf{n}}\cdot\hat{\mathbf{S}}_{\mathbf{n}+\delta}\tag{6.65}
\end{equation*}
在 $|\mathbf{q}\cdot\delta_{\max}|\ll1$ 时（$\delta_{\max}$ 是使 $\mathsf{J}(\delta)$
有任何可观贡献的最大矢量 $\delta$）有自旋波色散
\begin{equation*}
\hbar\omega(\mathbf{q}) = \hbar\omega(0) + S\sum_{\delta}\mathsf{J}(\delta)(\mathbf{q}\cdot\delta)^{2}\tag{6.66}
\end{equation*}

\exer{6.11} 式 6.23 中的态 $|q\rangle$ 表示未翻转自旋背景中的一个离域化的单个
翻转自旋。三维中它成为
\begin{equation*}
|\mathbf{q}\rangle = \frac{1}{\sqrt{N}}\sum_{\mathbf{j}}e^{i\mathbf{q}\cdot\mathbf{R}_{\mathbf{j}}}|\mathbf{j}\rangle.\tag{6.67}
\end{equation*}
翻转自旋位于格点 i 的概率为
\begin{equation*}
P = |\langle\mathbf{q}|\mathbf{i}\rangle|^{2}.\tag{6.68}
\end{equation*}
证明 $P=1/N$，从而翻转自旋以等概率分布。

横向自旋关联函数的算符为
\begin{equation*}
\hat{S}_{\mathbf{j}}^{x}\hat{S}_{\mathbf{j}'}^{x} + \hat{S}_{\mathbf{j}}^{y}\hat{S}_{\mathbf{j}'}^{y}.\tag{6.69}
\end{equation*}
证明
\begin{equation*}
\langle\mathbf{q}|\hat{S}_{\mathbf{j}}^{x}\hat{S}_{\mathbf{j}'}^{x}+\hat{S}_{\mathbf{j}}^{y}\hat{S}_{\mathbf{j}'}^{y}|\mathbf{q}\rangle = \frac{2S}{N}\cos(\mathbf{q}\cdot(\mathbf{R}_{\mathbf{j}}-\mathbf{R}_{\mathbf{j}'}))\tag{6.70}
\end{equation*}
（$j\neq j'$），并参照图 6.13 解释这一结果。

\exer{6.12} 我们对布洛赫壁宽的计算假定：壁内 N 个自旋中相邻自旋夹角处处为
$\pi/N$。这显然是个近似！下面是改进的计算。

设自旋作为垂直于壁面的坐标 $z$ 的函数由角 $\theta(z)$ 描述：$z\to-\infty$ 时
$\theta(z)\to0$，$z\to+\infty$ 时 $\theta(z)\to\pi$（对应 180° 壁）。能量为
\begin{equation*}
E = \int_{-\infty}^{\infty}\left[A\left(\frac{\partial\theta}{\partial z}\right)^{2} + f(\theta)\right]\mathrm{d}z.\tag{6.71}
\end{equation*}
A 是连续介质交换常数，$f(\theta)$ 描述各向异性。证明能量在
\begin{equation*}
\frac{\partial f}{\partial\theta} - 2A\frac{\partial^{2}\theta}{\partial z^{2}} = 0.\tag{6.72}
\end{equation*}
时极小。把该方程乘以 $\partial\theta/\partial z$ 并对 z 积分，证明
\begin{equation*}
f(\theta) = A\left(\frac{\partial\theta}{\partial z}\right)^{2}.\tag{6.73}
\end{equation*}
现在设
\begin{equation*}
f(\theta) = K\sin^{2}\theta\tag{6.74}
\end{equation*}
（K 为各向异性常数）。证明这意味着
\begin{equation*}
\cos\theta(z) = -\tanh\left(\frac{\pi z}{\delta}\right),\tag{6.75}
\end{equation*}
其中 $\delta=\pi\sqrt{A/K}$。画出 $\theta(z)$ 的草图并论证 $\delta$ 是畴壁宽度
的合适度量。再证明此情形下
\begin{equation*}
E = 4\sqrt{AK}\tag{6.76}
\end{equation*}

\exer{6.13} 考虑铁磁薄膜中的封闭畴结构（见图 6.30）：易轴竖直、难轴水平。长 180°
壁单位面积的能量为 $\sigma_{\mathrm{W}}$，各向异性能量密度为 K。忽略 90° 壁的
贡献（即设 $L\gg D$），证明
\begin{equation*}
D \approx \sqrt{\frac{2\sigma_{\mathrm{W}}L}{K}}.\tag{6.77}
\end{equation*}
取 $\sigma_{\mathrm{W}}=2\times10^{-3}$ J\,m$^{-2}$、$K=4\times10^{4}$ J\,m$^{-3}$、
$L=4$ mm，估算 D。

\begin{figure}
\centering
\includegraphics[width=0.4\textwidth]{fig6_30.pdf}
\caption{一个封闭畴。}
\end{figure}
